% ============================================================
% kbs_related_work.tex — KBS journal Related Work (KE voice) with a
% systematic comparison table. Section label: sec:rel (referenced by the
% Introduction draft, kbs_introduction.tex).
% Bibliography: kbs_ke_references.bib (new, verified) +
% paper/references.bib (conference refs carried over).
% ============================================================

\section{Related Work}
\label{sec:rel}

This section positions the paper in knowledge-engineering terms and then
compares prior work systematically. We organize the literature around the
stages of the knowledge-engineering cycle that this paper completes
---\textit{acquisition}, \textit{representation}, \textit{validation}, and
\textit{utilization}~\cite{studer_knowledge_engineering}---and place prior
work on LLM knowledge acquisition and metadata-aware retrieval within those
stages.

\subsection{Knowledge acquisition: from experts to LLMs}

The cost of acquiring structured knowledge has been the central bottleneck
of knowledge-based systems since the field's
origins~\cite{feigenbaum_knowledge_engineering}. Classical acquisition
relied on knowledge engineers eliciting expert knowledge, a process that is
slow, expensive, and hard to audit; principled knowledge-engineering
methodologies were developed to make each step
explicit~\cite{studer_knowledge_engineering}. Large language models
have changed the economics of the first step: they can propose structured
knowledge---ontology instances, graph triples, or metadata---from natural
language at negligible marginal cost~\cite{ciatto_llm_oracle_kbs,
pan_llm_kg_roadmap, xu_llm_gen_ie}. Recent surveys position LLMs as a
general acquisition channel for knowledge
graphs~\cite{ji_kg_survey, pan_llm_kg_roadmap} and for information
extraction~\cite{xu_llm_gen_ie}. What these works share is an emphasis on
\textit{generating} knowledge; validation against an independent gold
standard, and measurement of the knowledge's effect on a consuming system,
remain comparatively thin---precisely the stages the classical cycle makes
mandatory~\cite{studer_knowledge_engineering,
okeefe_kbs_validation}.

\subsection{Representation: thesaurus vs.\ taxonomy vs.\ ontology}

How acquired knowledge is represented determines what a consuming system can
do with it. Ontology engineering formalizes this choice, distinguishing
conceptual, representational, and encoding levels of an
ontology~\cite{gruber_portable_ontologies}. In the bibliographic domain,
two mature, professionally maintained representations coexist and are
directly comparable. \textit{Subject headings} (e.g., Library of Congress
Subject Headings) are an open, pre-coordinated thesaurus: a large,
human-curated vocabulary of heading strings that may combine facets (e.g.,
``Land tenure -- England -- History -- To 1500''). \textit{Classification}
(e.g., Dewey Decimal Classification) is a closed, hierarchical taxonomy of
numeric classes. The two differ in structure, in how new knowledge is
added, and---as we show---in how useful machine-acquired instances of each
are to a retrieval system. Prior metadata-aware retrieval work weights
structured fields but does not treat the \textit{type} of representation as
a variable~\cite{bm25f, prf_bm25}.

\subsection{Validation of machine-acquired knowledge}

The knowledge-validation literature established that agreement with a single
reference is not validation: expert systems must be checked for
completeness, consistency, and correctness against independent
standards~\cite{okeefe_kbs_validation}. Applied to LLMs, this standard is
rarely met. Studies of LLM subject cataloging report agreement with a
catalog record---one source---and stop
there~\cite{holstrom_llm_cataloging, characterising_ai_cataloguing,
ital_genai_cataloging, dobreski_chatbots_cataloging}. None measures an
inter-cataloger agreement ceiling (two professionals do not agree exactly on
every heading), and none checks proposed headings against the authority
file as part of the scoring protocol. The ``no free lunches'' finding for
LLM text post-correction---automation quality cannot be assumed and must be
measured against a proper baseline---applies with equal force to LLM
knowledge acquisition~\cite{ocr_no_free_lunches}.

\subsection{Utilization: metadata-aware and knowledge-based retrieval}

Retrieval systems that consume structured metadata are the natural test bed
for acquired knowledge. Fielded retrieval has long exploited weighted
fields~\cite{bm25f, prf_bm25}, and retrieval-augmented generation (RAG)
benchmarks standardize document-level evaluation but treat documents as
flat text~\cite{bergen}; hybrid semantic variants operate at the document
level as well~\cite{hysemrag}. Library-domain RAG systems adapt retrieval to
cultural-heritage and academic collections~\cite{folkrag,
lund_rag_libraries, historical_archives_rag, icadl_ephemeral} or to
indexing support~\cite{alima, koha_marc21_hitl}. These works establish
feasibility but leave the central question open: none measures whether
\textit{machine-acquired} metadata---as opposed to professionally created
metadata---is good enough to feed the retrieval system, let alone which
representation of it carries the effect and how the effect depends on how
much knowledge the base already has.

\subsection{Systematic comparison}
\label{sec:rel:table}

Table~\ref{tab:rel} contrasts the prior systems along the dimensions of the
knowledge-engineering cycle: acquisition channel, representation used or
produced, validation standard, whether downstream utility in a consuming
system is measured, and whether a released benchmark accompanies the work.
Entries are characterized from the papers' own abstracts and, where the
full text was available, their reported evaluation.

\begin{table*}[t]
\centering
\caption{Related systems positioned across the knowledge-engineering
cycle. $\checkmark$ = addressed; $\times$ = not addressed; $\circ$ =
partially/addressed incidentally. Rows are characterized from the cited
papers' reported scope; where a paper does not state a dimension we mark it
$\times$ rather than infer it.}
\label{tab:rel}
\begin{tabular}{lp{2.6cm}p{2.2cm}p{2.6cm}p{2.4cm}p{1.4cm}}
\toprule
Work & Acquisition channel & Representation & Validation standard &
Downstream utility measured & Released benchmark \\
\midrule
BERGEN~\cite{bergen} & $\times$ (uses existing docs) & flat text &
retrieval QA & retrieval nDCG & $\checkmark$ \\
HySemRAG~\cite{hysemrag} & $\times$ & flat text & QA-style &
literature-synthesis & $\times$ \\
FolkRAG~\cite{folkrag} & $\times$ & flat text + heritage metadata &
none stated & demonstration & $\times$ \\
Lund et al.~\cite{lund_rag_libraries} & $\times$ & library metadata &
none stated & feasibility analysis & $\times$ \\
ALIMA~\cite{alima} & LLM (indexing drafts) & GND thesaurus / DDC &
single-source (GND) & none stated & $\times$ \\
ChatGPT/Copilot
cataloging~\cite{holstrom_llm_cataloging} & LLM & LCSH/Sears, LCC &
single-source agreement & $\times$ & $\times$ \\
AI chatbot
cataloging~\cite{dobreski_chatbots_cataloging} & LLM & subject headings &
single-source agreement & $\times$ & $\times$ \\
GenAI cataloging
pilot~\cite{ital_genai_cataloging} & 4 LLMs & descriptive metadata
(RDA/MARC) & human review (small set) & $\times$ & $\times$ \\
Koha MARC21
HITL~\cite{koha_marc21_hitl} & LLM + human-in-loop & MARC 21 &
human review & $\times$ & $\times$ \\
Characterising AI
cataloguing~\cite{characterising_ai_cataloguing} & LLMs &
LCSH/classification & single-source agreement & $\times$ & $\times$ \\
LLM ontology
oracles~\cite{ciatto_llm_oracle_kbs} & LLM & ontology instances &
gold ontology & $\times$ (no retrieval) & $\times$ \\
\textbf{Ours} & \textbf{LLM} & \textbf{LCSH (thesaurus) + DDC (taxonomy)}
& \textbf{multi-level, authority-checked, agreement ceiling
\TODO{E2}} & \textbf{retrieval nDCG by representation and coverage}
& \textbf{$\checkmark$} \\
\bottomrule
\end{tabular}
\end{table*}

The comparison makes the gap explicit. On the acquisition side, LLM
cataloging studies validate against a single source
~\cite{holstrom_llm_cataloging, characterising_ai_cataloguing,
dobreski_chatbots_cataloging, ital_genai_cataloging} or propose
human-in-the-loop workflows~\cite{koha_marc21_hitl}, but none completes the
cycle: no multi-level rubric against independent gold, no authority-file
check in the scoring protocol, and no measurement of what the acquired
knowledge does to a retrieval system. On the utilization side, metadata-aware
retrieval work~\cite{bm25f, prf_bm25, bergen, hysemrag, folkrag,
lund_rag_libraries, alima} consumes metadata as given and never treats
machine acquisition as the variable. The one work closest to ours in method,
LLMs as oracles for ontology instantiation~\cite{ciatto_llm_oracle_kbs},
validates against gold ontology instances but does not measure downstream
utility in a consuming system. Ours is the only study, to our knowledge,
that carries LLM-acquired subject knowledge through all four stages---from
acquisition, through representation and validation, to measured utility in a
knowledge-based retrieval system---and reports the result as a function of
the representation and of knowledge-base completeness.
